\chapter{Statistika dan Peluang}\label{ch:g9:statproba}

Data dan kebetulan ada di mana-mana: hasil olahraga, ramalan cuaca,
permainan. Bab ini mengajarkan cara merangkum sederet bilangan dengan
\omterm{def:g9:statproba:indicators}{rata-rata}, \omterm{def:g9:statproba:indicators}{median} dan \omterm{def:g9:statproba:indicators}{jangkauannya}, serta cara menghitung \omterm{def:g9:statproba:probability}{peluang}
percobaan acak yang sederhana --- termasuk percobaan dua langkah, yang
ditangani dengan diagram pohon. Kedua pokok bahasannya dikembangkan
jauh lebih lanjut di buku Sekolah Menengah pada seri ini.

\section{Merangkum data}

\begin{definition}[Rata-rata, median, jangkauan]\label{def:g9:statproba:indicators}
Untuk sederet $N$ bilangan:
\begin{itemize}
  \item \emph{rata-rata}\index{rata-rata} adalah \omterm{def:g1:addition:def}{jumlah} semua nilainya
        dibagi $N$;
  \item \emph{median}\index{median} adalah nilai yang membelah deret
        terurutnya menjadi dua separuh yang sama besar: untuk $N$
        \omterm{def:g3:numbers:evenodd}{ganjil}, yaitu nilai tengahnya; untuk $N$ \omterm{def:g3:numbers:evenodd}{genap}, yaitu titik
        tengah kedua nilai tengahnya;
  \item \emph{jangkauan}\index{jangkauan} adalah nilai terbesar
        dikurangi nilai terkecil.
\end{itemize}
\end{definition}

\begin{example}\label{ex:g9:statproba:indicators}
Nilai seorang murid: $8,\ 12,\ 9,\ 15,\ 11$.
\begin{enumerate}
  \item \omterm{def:g9:statproba:indicators}{Rata-rata}: $\dfrac{8 + 12 + 9 + 15 + 11}{5} = \dfrac{55}{5} = 11$.
  \item \omterm{def:g9:statproba:indicators}{Median}: urutkan deretnya: $8, 9, 11, 12, 15$; nilai tengahnya
        (yang ke-3) adalah $11$.
  \item \omterm{def:g9:statproba:indicators}{Jangkauan}: $15 - 8 = 7$.
\end{enumerate}
Dengan nilai keenam $17$: terurut menjadi $8, 9, 11, 12, 15, 17$,
\omterm{def:g9:statproba:indicators}{mediannya} menjadi titik tengah $11$ dan $12$, yaitu $11.5$, dan
\omterm{def:g9:statproba:indicators}{rata-ratanya} menjadi $\frac{72}{6} = 12$.
\end{example}

\begin{example}[Rata-rata berbobot]\label{ex:g9:statproba:weighted}
Ukuran sepatu yang terjual dalam sehari:
\begin{center}
\renewcommand{\arraystretch}{1.2}
\begin{tabular}{c|ccccc}
ukuran & $38$ & $39$ & $40$ & $41$ & $42$ \\
\hline
\omterm{def:g7:stats:frequency}{frekuensi} & $3$ & $7$ & $6$ & $3$ & $1$ \\
\end{tabular}
\end{center}
Ukuran \omterm{def:g9:statproba:indicators}{rata-ratanya} adalah
\[
\frac{3 \times 38 + 7 \times 39 + 6 \times 40 + 3 \times 41 + 1 \times 42}
{3 + 7 + 6 + 3 + 1}
= \frac{114 + 273 + 240 + 123 + 42}{20}
= \frac{792}{20} = 39.6 .
\]
\end{example}

\begin{remark}
\omterm{def:g9:statproba:indicators}{Rata-rata} dan \omterm{def:g9:statproba:indicators}{median} dapat berbeda jauh. Pada deret
$10, 10, 10, 10, 60$ \omterm{def:g9:statproba:indicators}{rata-ratanya} $20$ (yang tertarik naik oleh nilai
$60$) sedangkan \omterm{def:g9:statproba:indicators}{mediannya} $10$. Tanyakanlah selalu penunjuk yang mana
yang mewakili keadaannya dengan lebih baik.
\end{remark}

\section{Peluang}

\begin{definition}[Peluang]\label{def:g9:statproba:probability}
Sebuah percobaan acak punya beberapa \emph{hasil} yang mungkin; sebuah
\emph{kejadian} adalah himpunan hasil. Tiap hasilnya memperoleh sebuah
\emph{peluang}\index{peluang}: yaitu bilangan antara $0$ dan $1$,
dengan semua peluang hasilnya berjumlah $1$; peluang $\P(A)$ sebuah
kejadian $A$ adalah \omterm{def:g1:addition:def}{jumlah} peluang hasilnya. Ketika seluruh $n$
hasilnya sama-sama mungkin,
\[
\P(A) = \frac{\text{banyaknya hasil yang menguntungkan}}{n}.
\]
\end{definition}

\begin{example}\label{ex:g9:statproba:spinner}
Sebuah roda terbagi menjadi $8$ juring yang sama: $3$ merah, $3$ biru,
$2$ hijau. Tiap juringnya berpeluang $\frac18$, jadi
\[
\P(\text{merah}) = \frac38, \qquad
\P(\text{hijau}) = \frac28 = \frac14, \qquad
\P(\text{bukan hijau}) = 1 - \frac14 = \frac34 .
\]
\end{example}

\begin{proposition}[Komplemen]\label{prop:g9:statproba:complement}
Untuk setiap kejadian $A$, \emph{kejadian sebaliknya} $\bar A$ (``$A$
tidak terjadi'') memenuhi
\[
\P(\bar A) = 1 - \P(A).
\]
\end{proposition}

\begin{proof}
Setiap hasilnya berada di tepat satu dari $A$ dan $\bar A$, jadi
$\P(A) + \P(\bar A)$ menjumlahkan \omterm{def:g9:statproba:probability}{peluang} semua hasilnya, yang bernilai
$1$.
\end{proof}

\section{Percobaan dua langkah}

\begin{method}[Diagram pohon]\label{met:g9:statproba:tree}
Untuk percobaan dalam dua langkah:
\begin{enumerate}
  \item gambarlah satu cabang untuk tiap hasil langkah pertamanya,
        lalu lanjutkan tiap cabangnya dengan hasil langkah keduanya,
        sambil menulis tiap \omterm{def:g9:statproba:probability}{peluangnya} pada cabangnya;
  \item kalikan \omterm{def:g9:statproba:probability}{peluang} di sepanjang satu jalur untuk memperoleh
        \omterm{def:g9:statproba:probability}{peluang} jalur itu;
  \item jumlahkan \omterm{def:g9:statproba:probability}{peluang} semua jalur yang membentuk sebuah kejadian.
\end{enumerate}
\end{method}

\begin{example}\label{ex:g9:statproba:tree}
Sebuah kantong memuat $2$ keping merah dan $3$ keping hitam. Ambillah
satu keping, catat warnanya, \emph{kembalikan}, lalu ambil lagi. Tiap
pengambilannya memberi merah dengan \omterm{def:g9:statproba:probability}{peluang} $\frac25$, dan hitam
dengan \omterm{def:g9:statproba:probability}{peluang} $\frac35$.

\begin{omfigure}
\begin{tikzpicture}[scale=0.95]
  \coordinate (root) at (0,0);
  \coordinate (R) at (2.2,1.0);
  \coordinate (B) at (2.2,-1.0);
  \coordinate (RR) at (4.7,1.55);
  \coordinate (RB) at (4.7,0.45);
  \coordinate (BR) at (4.7,-0.45);
  \coordinate (BB) at (4.7,-1.55);
  \draw (root) -- (R) node[midway, above, font=\small, sloped] {$\frac25$};
  \draw (root) -- (B) node[midway, below, font=\small, sloped] {$\frac35$};
  \draw (R) -- (RR) node[midway, above, font=\small, sloped] {$\frac25$};
  \draw (R) -- (RB) node[midway, below, font=\small, sloped] {$\frac35$};
  \draw (B) -- (BR) node[midway, above, font=\small, sloped] {$\frac25$};
  \draw (B) -- (BB) node[midway, below, font=\small, sloped] {$\frac35$};
  \node[omThm, font=\small, above left=-2pt] at (R) {M};
  \node[font=\small, below left=-2pt] at (B) {H};
  \node[right, font=\small] at (RR)
    {MM:\ $\frac25 \times \frac25 = \frac{4}{25}$};
  \node[right, font=\small] at (RB)
    {MH:\ $\frac25 \times \frac35 = \frac{6}{25}$};
  \node[right, font=\small] at (BR)
    {HM:\ $\frac35 \times \frac25 = \frac{6}{25}$};
  \node[right, font=\small] at (BB)
    {HH:\ $\frac35 \times \frac35 = \frac{9}{25}$};
\end{tikzpicture}

{\small Pohon kedua pengambilan dengan pengembalian: kalikan sepanjang
cabangnya, lalu periksalah bahwa keempat jalurnya berjumlah $1$.}
\end{omfigure}

\omterm{def:g9:statproba:probability}{Peluang} dua keping yang warnanya sama:
$\frac{4}{25} + \frac{9}{25} = \frac{13}{25}$. \omterm{def:g9:statproba:probability}{Peluang} sedikitnya
satu merah: $1 - \P(\text{HH}) = 1 - \frac{9}{25} = \frac{16}{25}$.
\end{example}

\begin{remark}[Frekuensi relatif mendekati peluang]
Melempar dadu yang setimbang $6000$ kali memberi \emph{kira-kira}
$1000$ angka enam --- bukan persis. Ketika banyaknya pengulangan
bertambah, \omterm{def:g7:stats:frequency}{frekuensi relatif} yang teramati menjadi makin dekat dengan
\omterm{def:g9:statproba:probability}{peluangnya}: itulah yang membuat \omterm{def:g9:statproba:probability}{peluang} menjadi alat yang tepat untuk
memodelkan percobaan berulang (yaitu kisah yang dikembangkan di buku
Sekolah Menengah pada seri ini).
\end{remark}

\section{Latihan}

\begin{exercise}[$\star$]\label{exo:g9:statproba:1}
Hitunglah \omterm{def:g9:statproba:indicators}{rata-rata}, \omterm{def:g9:statproba:indicators}{median} dan \omterm{def:g9:statproba:indicators}{jangkauan} deret
$14,\ 9,\ 17,\ 9,\ 12,\ 11,\ 12$.
\end{exercise}

\begin{exercise}[$\star$]\label{exo:g9:statproba:2}
Suhu pada tengah hari selama sepekan adalah $19,\ 21,\ 24,\ 18,\ 22,\
25,\ 19$ (derajat). Hitunglah \omterm{def:g9:statproba:indicators}{rata-ratanya} (sampai persepuluhan) dan
\omterm{def:g9:statproba:indicators}{mediannya}.
\end{exercise}

\begin{exercise}[$\star$]\label{exo:g9:statproba:3}
Banyaknya gol yang dicetak sebuah tim pada $10$ pertandingan:
\begin{center}
\renewcommand{\arraystretch}{1.2}
\begin{tabular}{c|cccc}
gol & $0$ & $1$ & $2$ & $3$ \\
\hline
pertandingan & $2$ & $4$ & $3$ & $1$ \\
\end{tabular}
\end{center}
Hitunglah \omterm{def:g9:statproba:indicators}{rata-rata} banyaknya gol tiap pertandingan, dan \omterm{def:g9:statproba:indicators}{mediannya}.
\end{exercise}

\begin{exercise}[$\star$]\label{exo:g9:statproba:4}
Sebuah dadu yang setimbang dilempar sekali. Hitunglah \omterm{def:g9:statproba:probability}{peluang}:
``memperoleh $3$''; ``memperoleh \omterm{def:g3:numbers:evenodd}{bilangan genap}''; ``memperoleh
sedikitnya $3$''; ``tidak memperoleh $6$''.
\end{exercise}

\begin{exercise}[$\star$]\label{exo:g9:statproba:5}
Sebuah kotak berisi $12$ bola: $5$ putih, $4$ merah dan $3$ hijau;
satu bola diambil secara acak. Hitunglah $\P(\text{putih})$,
$\P(\text{merah atau hijau})$, dan $\P(\text{bukan putih})$. Apa yang
kamu perhatikan tentang kedua yang terakhir?
\end{exercise}

\begin{exercise}[$\star\star$]\label{exo:g9:statproba:6}
Sebuah koin dilempar, lalu sebuah dadu yang setimbang dilempar.
Gambarlah pohonnya (atau bilanglah hasilnya) lalu hitunglah \omterm{def:g9:statproba:probability}{peluang}
memperoleh sisi gambar \emph{dan} angka $6$; lalu \omterm{def:g9:statproba:probability}{peluang} memperoleh
sisi gambar \emph{atau} angka $6$ (atau keduanya).
\end{exercise}

\begin{exercise}[$\star\star$]\label{exo:g9:statproba:7}
Pada kantong \cref{ex:g9:statproba:tree} ($2$ merah, $3$ hitam), kedua
pengambilannya kini dilakukan \emph{tanpa} pengembalian. Gambarlah
pohon yang baru (hati-hati: \omterm{def:g9:statproba:probability}{peluang} tingkat keduanya berubah) lalu
hitunglah \omterm{def:g9:statproba:probability}{peluang} mengambil dua keping hitam, lalu \omterm{def:g9:statproba:probability}{peluang} mengambil
dua keping yang warnanya berbeda.
\end{exercise}

\begin{exercise}[$\star\star$]\label{exo:g9:statproba:8}
Sesudah $9$ pertandingan, seorang pemain bola basket \omterm{def:g9:statproba:indicators}{rata-rata}
mencetak $14$ angka tiap pertandingan.
\begin{enumerate}
  \item Berapa angka yang sudah dicetaknya seluruhnya?
  \item Berapa angka yang harus dicetaknya pada pertandingan ke-$10$
        agar \omterm{def:g9:statproba:indicators}{rata-ratanya} $15$?
\end{enumerate}
\end{exercise}

\begin{exercise}[$\star\star\star$]\label{exo:g9:statproba:9}
Sebuah permainan: lemparkan dua dadu yang setimbang; kamu menang kalau
kedua hasilnya sama (yaitu ``sebuah dobel'').
\begin{enumerate}
  \item Hitunglah \omterm{def:g9:statproba:probability}{peluang} menang.
  \item Kamu bermain dua kali (dengan permainan yang saling bebas).
        Dengan memakai pohon, hitunglah \omterm{def:g9:statproba:probability}{peluang} menang sedikitnya
        sekali.
\end{enumerate}
\end{exercise}

\section{Soal: Taruhan Chevalier yang membangkrutkan}

\begin{problem}[{Soal akhir pekan --- perselisihan judi tahun 1654
yang menciptakan teori peluang, dan kebetulan ulang tahun yang
mengelabui setiap kelas}]
\label{pb:g9:statproba:1}
Pada tahun 1654 seorang bangsawan Prancis yang gemar berjudi, yaitu
Chevalier de Méré, menjadi kaya lewat satu taruhan dadu lalu mulai
kalah pada taruhan lain yang dipercayanya setara. Karena bingung, ia
bertanya kepada ahli matematika Blaise Pascal; Pascal menulis surat
kepada Pierre de Fermat; dan pertukaran surat mereka mendirikan teori
\omterm{def:g9:statproba:probability}{peluang}. Soal ini memainkan ulang seluruh peristiwanya dengan alat bab
ini --- hasil yang sama-sama mungkin, pohon dan aturan komplemennya
(\cref{prop:g9:statproba:complement}) --- lalu berakhir dengan
kebetulan yang paling berlawanan dengan gerak hati di antara semuanya.

\medskip
\noindent\textbf{Bagian I --- Dadu, yang dibilang secara jujur.}
\begin{enumerate}
  \item Lemparkan dadu yang setimbang dua kali. Dengan memakai pohon
        (atau tabel $6 \times 6$), hitunglah \omterm{def:g9:statproba:probability}{peluang} \emph{tidak} ada
        angka enam pada kedua lemparannya, lalu simpulkan \omterm{def:g9:statproba:probability}{peluang}
        sedikitnya satu angka enam.
  \item Seorang teman sekelas berdalih: ``satu \omterm{def:g9:statproba:probability}{peluang} dari enam tiap
        lemparan, dua kali, jadi $\frac16 + \frac16 = \frac13$.''
        Bandingkan dengan pertanyaan 1 lalu tunjukkan hasil persis
        mana yang dibilang penjumlahannya
        dua kali.
  \item Melempar dua dadu lalu menjumlahkannya: hitunglah \omterm{def:g9:statproba:probability}{peluang}
        \omterm{def:g1:addition:def}{jumlahnya} $7$ dan \omterm{def:g9:statproba:probability}{peluang} \omterm{def:g1:addition:def}{jumlahnya} $2$. Mengapa $7$ menjadi
        kegemaran para penjudi?
  \item Perselisihan lama: dengan dua dadu, \omterm{def:g1:addition:def}{jumlah} $9$ atau \omterm{def:g1:addition:def}{jumlah}
        $10$ yang lebih mungkin? Bilanglah hasilnya lalu putuskan.
  \item Perumumlah pertanyaan 1: jelaskan mengapa \omterm{def:g9:statproba:probability}{peluang} sedikitnya
        satu angka enam dalam $n$ lemparan adalah
        \[
        1 - \left(\frac56\right)^{\!n} ,
        \]
        karena cabang tanpa angka enam pada pohonnya mengalikan diri
        dari tingkat ke tingkat.
\end{enumerate}

\medskip
\noindent\textbf{Bagian II --- Kedua taruhan Chevalier.}
\begin{enumerate}[resume]
  \item Taruhan pertama, dengan uang seimbang: \emph{sedikitnya satu
        angka enam dalam empat lemparan sebuah dadu}. Hitunglah
        $\left(\frac56\right)^4$ sebagai \omterm{def:g9:fractions:fraction}{pecahan} dan sebagai \omterm{def:g6:decimals:places}{bilangan
        desimal}, lalu \omterm{def:g9:statproba:probability}{peluang} menang Chevalier itu. Apakah taruhannya
        bagus?
  \item Penalaran Chevalier sendiri adalah: ``empat lemparan pada
        $\frac16$ masing-masing: jadi $\frac46$ \omterm{def:g9:statproba:probability}{peluang}.'' Penalaran
        itu memberi jawaban yang hampir benar di sini --- tetapi
        doronglah ke $7$ lemparan: kemustahilan apa yang
        dihasilkannya, dan kekeliruan pertanyaan 2 yang mana yang
        diulanginya?
  \item Taruhan kedua, yang dipercayanya setara lewat \omterm{def:g9:linfunc:linear}{perbandingan
        senilai} (``$24$ lemparan pada $\frac1{36}$ masing-masing:
        lagi-lagi $\frac{24}{36} = \frac46$''): \emph{sedikitnya satu
        dobel enam dalam $24$ lemparan dua dadu}. Hitunglah \omterm{def:g9:statproba:probability}{peluang}
        menang yang sebenarnya $1 - \left(\frac{35}{36}\right)^{24}$
        (dengan kalkulator). Mengapa Chevalier perlahan bangkrut?
  \item Berapa lemparan dua dadu yang akan memulihkan keunggulannya?
        Ujilah $n = 25$ dengan kalkulator lalu simpulkan.
  \item Nyatakan pesan moralnya dalam dua kalimat: apa yang keliru,
        secara umum, dengan mengalikan sebuah \omterm{def:g9:statproba:probability}{peluang} dengan
        banyaknya percobaan --- dan alat yang benar yang mana yang
        menggantikan godaan itu?
\end{enumerate}

\medskip
\noindent\textbf{Bagian III --- \omterm{def:g9:statproba:indicators}{Rata-rata}, \omterm{def:g9:statproba:indicators}{median}, ulang tahun.}
\begin{enumerate}[resume]
  \item Sebuah perusahaan kecil membayar sembilan pegawainya
        $2\,000$ euro tiap bulan dan direkturnya $20\,000$. Hitunglah
        gaji \omterm{def:g9:statproba:indicators}{rata-rata} dan gaji \omterm{def:g9:statproba:indicators}{mediannya}
        (\cref{def:g9:statproba:indicators}). Bilangan yang mana yang
        sebaiknya dikutip iklan lowongannya secara jujur?
  \item Susunlah daftar lima nilai ujian dengan \omterm{def:g9:statproba:indicators}{rata-rata} $12$ dan
        \omterm{def:g9:statproba:indicators}{median} $15$. Apa yang dikatakan pasangan itu (\omterm{def:g9:statproba:indicators}{rata-rata} di
        bawah \omterm{def:g9:statproba:indicators}{median}) tentang bentuk sebaran nilainya?
  \item Soal ulang tahun: pada kelas berisi $23$ murid, berapa
        \omterm{def:g9:statproba:probability}{peluang} bahwa sedikitnya dua murid berulang tahun sama?
        Susunlah komplemennya (semua $23$ ulang tahunnya berbeda):
        \[
        \frac{365}{365} \times \frac{364}{365} \times
        \frac{363}{365} \times \dots \times \frac{343}{365} .
        \]
        Hitunglah \omterm{def:g2:mult:def}{hasil kali} ketiga faktor pertamanya, uraikan
        bagaimana faktornya berkembang, lalu --- karena \omterm{def:g2:mult:def}{hasil kali}
        seluruhnya kira-kira $0.493$ --- jawablah pertanyaannya.
        Kebanyakan orang menduga ``sangat tidak mungkin'': apa kata
        matematikanya?
  \item Gerak hati yang diperbaiki: berapa \emph{pasang} murid yang
        dapat berulang tahun sama pada kelas berisi $23$ murid
        (yaitu bilangan jabat tangan pada \cref{pb:g6:wholes:1})?
        Jelaskan dalam satu kalimat mengapa bilangan itu, bukan $23$
        itu sendiri, yang menggerakkan kejutannya.
  \item Penutupnya: uraikan sebuah percobaan untuk memeriksa salah
        satu hasilnya --- taruhan de Méré dengan dadu sungguhan, atau
        soal ulang tahun pada beberapa kelas --- lalu nyatakan dengan
        hati-hati apa yang kamu harapkan dari \omterm{def:g7:stats:frequency}{frekuensi relatif} yang
        teramati ketika banyaknya pengulangan bertambah. (Harapan itu
        punya nama, yaitu hukum bilangan besar; buktinya menunggu di
        buku Sekolah Menengah.)

\end{enumerate}
\end{problem}
